% Contoh prediksi 2 hari data uji (MLP INT8 = chip) vs aktual.
\documentclass[border=4pt]{standalone}
% Gaya bersama gambar deck Demo Day (dipakai semua gambar/*.tex standalone).
\usepackage{amssymb}
\usepackage{fontspec}
\setmainfont{DejaVu Sans Condensed}
\setsansfont{DejaVu Sans Condensed}
\newfontfamily\judulfont{DejaVu Sans}[BoldFont={DejaVu Sans Bold}, Ligatures=TeX]
\usepackage{tikz,pgfplots}
\pgfplotsset{compat=1.18, every axis/.append style={/pgf/number format/assume math mode=true, /pgf/number format/use comma, /pgf/number format/1000 sep={.}}}
\usetikzlibrary{positioning,arrows.meta,calc,shapes.geometric,fit,backgrounds}
\definecolor{gelap}{HTML}{15242B}
\definecolor{teal}{HTML}{285052}
\definecolor{amber}{HTML}{E8A317}
\definecolor{muda}{HTML}{DCE8E8}
\definecolor{abu}{HTML}{8A9499}
\definecolor{teks}{HTML}{1E2B30}
\definecolor{hijau}{HTML}{2F7D5B}
\definecolor{merah}{HTML}{A33A2A}
\color{teks}
\tikzset{
  ubin/.style={fill=muda, rounded corners=5pt, align=center, inner sep=6pt},
  inti/.style={ubin, fill=gelap, text=white},
  panah/.style={-{Stealth[length=3mm]}, line width=1.3pt, draw=amber},
  panahabu/.style={-{Stealth[length=2.5mm]}, line width=1pt, draw=abu},
  judulubin/.style={font=\judulfont\bfseries},
}
% Dibuat otomatis oleh buat_aset.py dari ml/results/energi_harian.json; jangan diedit manual.
\newcommand{\energiTurbin}{270}
\newcommand{\energiQa}{187}
\newcommand{\energiQb}{377}
\newcommand{\hariEnergi}{140}
\newcommand{\raspiWh}{71}
\newcommand{\chipWh}{4,1}
\newcommand{\raspiPersen}{26}
\newcommand{\chipPersen}{1,5}
\newcommand{\raspiPersenA}{19}
\newcommand{\raspiPersenB}{38}
\newcommand{\raspiPersenNum}{26.3}
\newcommand{\chipPersenNum}{1.5}
\newcommand{\hematKwh}{24}
\newcommand{\kaliHemat}{17}
\newcommand{\jRaspi}{182}
\newcommand{\jChip}{10,4}
\newcommand{\periodeEnergi}{2026-03-31--2026-10-20}

\AtBeginDocument{\hyphenpenalty=10000 \exhyphenpenalty=10000}

\usepgfplotslibrary{fillbetween}
\begin{document}
\begin{tikzpicture}[font=\fontsize{9}{11}\selectfont]
\begin{axis}[width=13cm, height=5cm, ymin=0, ymax=135, xmin=0, xmax=48, xtick={0,6,...,48}, ylabel={W, rata-rata 1 jam ke depan}, xlabel={jam},
  axis line style={abu}, tick style={abu}, legend style={draw=none, fill=none, at={(0.99,0.99)}, anchor=north east, legend columns=3}]
  \addplot[name path=lo, draw=none, forget plot] table[col sep=comma, x=jam, y=p10] {../../../data/contoh_prediksi.csv};
  \addplot[name path=hi, draw=none, forget plot] table[col sep=comma, x=jam, y=p90] {../../../data/contoh_prediksi.csv};
  \addplot[muda, forget plot] fill between[of=lo and hi];
  \addplot[gelap, line width=1.2pt] table[col sep=comma, x=jam, y=aktual] {../../../data/contoh_prediksi.csv};
  \addplot[amber, line width=1.6pt] table[col sep=comma, x=jam, y=int8] {../../../data/contoh_prediksi.csv};
  \legend{aktual, prediksi chip}
\end{axis}
\end{tikzpicture}
\end{document}
